\begin{abstract}
物体領域の検出とカテゴリのアノテーションを分離し，\rev{WordNetとVisual Genome（VG）を組み合わせた語彙の木}をたどってカテゴリの粒度を変更する，再学習を必要としない物体アノテーションのモデルを示す．
本モデルは，検出器にOWLv2のクラス非依存の物体らしさ，分類器に注目領域をマスクで与えるAlpha-CLIPを用い，\rev{語彙の木の全ての語}の確率を語彙の木の部分木ごとに合計して，確信度に応じた深さの語を出力する．
検出器と分類器をそれぞれ置き換えた場合と比較し，各部品を選ぶ理由を示す．
COCO 2017 Validationの500枚による評価では，OWLv2は1枚あたり50個のボックスでCOCOの物体の87.5\%を囲み，DETRの73.8\%とYOLO11xの70.2\%を上回った．
同じ画像のLVISの注釈によるCOCO外の物体の再現率は，OWLv2の73.4\%に対してYOLO11xは10.0\%，DETRは12.4\%であり，COCOで学習した検出器はCOCO外の物体をほとんど囲めなかった．
一方，COCOの80クラスによるF値は，検出器をYOLO11xのボックスに替えた構成が0.535，DETRが0.486，OWLv2が0.411であり，COCOの物体に限ればCOCOで学習した検出器のボックスが有利であった．
分類器は，いずれの検出器でもAlpha-CLIPのF値がCLIPより0.03〜0.05高く，COCOのクラスより細かい語を正しく出力した割合もCLIPの約2倍であった．
GTのボックスに対する命名では，部分木の確率の合計により，COCOのクラスのみで命名する場合と比べて階層的F値が0.266から0.639に向上した．
\end{abstract}

\noindent
{\bf キーワード：}物体アノテーション，粒度可変アノテーション，WordNet，Visual Genome，OWLv2，Alpha-CLIP，YOLO，DETR，オープンボキャブラリ
